\documentclass[conference]{IEEEtran}
\usepackage{cite}
\usepackage{amsmath,amssymb,amsfonts}
\usepackage{algorithmic}
\usepackage{graphicx}
\usepackage{textcomp}
\usepackage{xcolor}
\usepackage{booktabs}
\usepackage{multirow}
\usepackage{hyperref}
\usepackage{float}
\usepackage{subcaption}
\usepackage{algorithm}
\usepackage{placeins}
\definecolor{pendingcolor}{RGB}{200, 100, 50}
\newcommand{\pending}[1]{\textcolor{pendingcolor}{\textit{[#1]}}}

\begin{document}

\title{Graph Structure Beats Tabular Boosting for Money-Laundering Detection: An Inductive GraphSAGE Study on SAML-D}

\author{
\IEEEauthorblockN{Prit Kudale\IEEEauthorrefmark{1}, Naman Dwivedi\IEEEauthorrefmark{2}, Raj Dandekar\IEEEauthorrefmark{2}, Rajat Dandekar\IEEEauthorrefmark{2}, Sreedath Panat\IEEEauthorrefmark{2}}
\IEEEauthorblockA{\IEEEauthorrefmark{1}prit.kudale@gmail.com \quad \IEEEauthorrefmark{2}Vizuara AI Labs, hello@vizuara.com}
}

\maketitle

\begin{abstract}
Money laundering is a relational, multi-hop phenomenon, yet most transaction-monitoring systems score
transactions independently. Whether graph learning actually helps, and how to train it under extreme class
imbalance, remains unsettled, and no controlled inductive comparison exists on the recent SAML-D benchmark.
We model SAML-D (9.5M transactions, 0.104\% laundering) as a directed graph in which nodes are accounts and
edges are transactions, undersample the majority class to a 247{,}247-edge subset (3.99\% positive), and
train an inductive GraphSAGE-mean encoder with learned categorical embeddings and an edge-MLP head,
evaluated on a stratified \emph{disjoint-node} split in which test accounts are unseen during training. We
benchmark against XGBoost on identical features, take recall (constrained to $\geq 0.70$) as the primary,
regulator-relevant metric, and compare focal loss against weighted cross-entropy. GraphSAGE decisively
outperforms XGBoost: test PR-AUC $0.741 \pm 0.013$ versus $0.288 \pm 0.006$ and roughly $6\times$ higher
precision ($0.564$ versus $0.089$) at recall $\geq 0.70$ (Welch $t$-test $p < 0.001$, three seeds). Focal
loss, however, is statistically on par with weighted cross-entropy (PR-AUC $0.736$ versus $0.741$,
$p = 0.63$). Relational structure, not the loss function, is the decisive lever for anti-money-laundering
on SAML-D, and the entire study runs in under one dollar on a single L4 GPU.
\end{abstract}

\begin{IEEEkeywords}
anti-money laundering, graph neural networks, GraphSAGE, class imbalance, focal loss, fraud detection
\end{IEEEkeywords}

\section{Introduction}
Money laundering moves illicit funds through multi-hop structures: a source account fans out to many
``mule'' accounts (smurfing), many accounts funnel into one (structuring), or funds cycle through a chain
of intermediaries (layering). These signatures live in the topology of the transaction network and are
largely invisible to models that score each transaction in isolation. Figure~\ref{fig:typologies}
illustrates the canonical typologies as graph motifs. Because a missed laundering transaction (a false
negative) exposes a financial institution to regulatory penalties and sanctions, while a false positive is
a recoverable investigation, recall is the operationally decisive metric for anti-money-laundering (AML).

Graph neural networks (GNNs) are a natural fit, yet two questions remain open in practice. First, does a
GNN actually beat a strong tabular model once both are given the same features? The evidence is mixed: on
the Elliptic Bitcoin dataset a Random Forest beat a graph convolutional network~\cite{weber2019aml}, and on
IBM AMLworld gradient-boosted trees were near-par with message-passing GNNs~\cite{altman2023amlworld}.
Second, how should training handle the extreme imbalance of real AML data (here, $0.104\%$ positives)?
Focal loss~\cite{lin2017focal} was designed for exactly this regime in vision, but whether it beats simple
class-weighted cross-entropy for graph fraud is untested. The recent SAML-D benchmark~\cite{oztas2023samld}
provides a realistic, fully labelled testbed, but its own paper evaluates only classical tabular models
(no GNN and no gradient boosting) and explicitly defers ``more complex ML algorithms'' to future work.

We close this gap with a controlled, budget-bounded study (Fig.~\ref{fig:overview}). Our
\textbf{contributions} are:
\begin{enumerate}
\item The first controlled, inductive, \emph{disjoint-node} GraphSAGE-versus-XGBoost edge-classification
benchmark on SAML-D, with recall-constrained ($\geq 0.70$) post-hoc threshold tuning and multi-seed error bars.
\item A decisive, statistically significant result that \textbf{graph structure helps}: GraphSAGE reaches
PR-AUC $0.74$ versus XGBoost's $0.29$ and roughly $6\times$ the precision at recall $\geq 0.70$
($p < 0.001$).
\item A controlled \textbf{focal-versus-weighted-cross-entropy} ablation showing the two are statistically
\emph{on par} for this task ($p = 0.63$); focal loss's vision-era advantage does not robustly transfer.
\item A \textbf{reproducible, low-cost protocol} (under one dollar on a single L4 GPU, $15$ training runs)
with released code and a $16$-run result ledger.
\end{enumerate}

\begin{figure*}[t]
\centering
\includegraphics[width=0.96\textwidth]{figures/output/fig_overview_pb.png}
\caption{End-to-end pipeline. SAML-D transactions are modelled as a directed account graph; an inductive
GraphSAGE encoder produces account embeddings; an edge-MLP over the sender and receiver embeddings plus
learned categorical features predicts $P(\text{is\_laundering})$. Extreme imbalance is handled by majority
undersampling and a recall-first operating point, and generalization is measured on a disjoint-node split.}
\label{fig:overview}
\end{figure*}

\section{Background and Related Work}
\subsection{AML and graph-based fraud detection}
SAML-D~\cite{oztas2023samld} is a synthetic transaction-monitoring dataset of $9{,}504{,}852$ transactions
with $28$ typologies and a $0.104\%$ laundering rate. The Elliptic study~\cite{weber2019aml} contributed a
Bitcoin AML graph and the cautionary finding that a Random Forest can beat a GCN. IBM
AMLworld~\cite{altman2023amlworld} frames AML as edge classification and reports gradient-boosted trees
near-par with vanilla GNNs. Recent surveys~\cite{cheng2024review} identify class imbalance as a central
open problem for GNN-based fraud detection.

\subsection{Inductive graph neural networks}
GraphSAGE~\cite{hamilton2017graphsage} learns aggregator functions (mean, LSTM, pooling) that embed a node
from its sampled neighborhood, enabling \emph{inductive} inference on unseen nodes, unlike the transductive
GCN~\cite{kipf2017gcn}. Attention-based GAT~\cite{velickovic2018gat} is inductive-capable but is excluded
here (see Section~\ref{sec:setup}). Specialized fraud GNNs such as CARE-GNN~\cite{dou2020caregnn} and
PC-GNN~\cite{liu2021pcgnn} add camouflage-resistant neighbor selection and label-balanced sampling; we keep
the encoder deliberately simple to respect the compute budget.

\subsection{Class imbalance}
Standard remedies are majority undersampling, class-weighted loss, and focal loss~\cite{lin2017focal},
combined with post-hoc threshold tuning. Under extreme imbalance, PR-AUC is more informative than ROC-AUC
because the large true-negative pool inflates specificity~\cite{davis2006pr}. Our tabular baseline is
XGBoost~\cite{chen2016xgboost}; the implementation uses PyTorch Geometric~\cite{fey2019pyg}.

\section{Methodology}
\label{sec:method}
\subsection{Problem formulation}
We build a directed graph $G = (V, E)$ where $V$ is the set of accounts and each transaction is an edge
$e = (s, r) \in E$ with features $x_e$. The task is binary edge classification of $\text{is\_laundering}$.
The GraphSAGE mean aggregator computes, at layer $k$,
\begin{equation}
h_v^{k} = \sigma\!\left( W^{k} \cdot \left[\, h_v^{k-1} \;\Vert\; \mathrm{mean}_{u \in \mathcal{N}(v)}\, h_u^{k-1} \,\right] \right),
\end{equation}
where $\Vert$ denotes concatenation and $\mathcal{N}(v)$ the neighbors of $v$. For a transaction edge we
score
\begin{equation}
\hat{y}_e = \mathrm{MLP}\!\left( \left[\, h_s \;\Vert\; h_r \;\Vert\; E_{\mathrm{cat}}(c_e) \;\Vert\; x^{\mathrm{num}}_e \,\right] \right),
\end{equation}
where $E_{\mathrm{cat}}$ are learned embeddings of the five categorical columns and $x^{\mathrm{num}}_e$ the
numeric (log-amount and temporal) features. The encoder and head are shown in
Fig.~\ref{fig:architecture}.

\subsection{Imbalance handling}
We undersample the majority class, keeping all laundering edges and a $2.5\%$ sample of normal edges, which
raises the positive rate from $0.104\%$ to $3.99\%$. We compare weighted binary cross-entropy,
\begin{equation}
\mathcal{L}_{\mathrm{wBCE}} = -\, w^{+} y \log \hat{y} - (1-y)\log(1-\hat{y}),
\end{equation}
with $w^{+}$ the negative-to-positive ratio, against focal loss
\begin{equation}
\mathrm{FL}(p_t) = -\, \alpha_t\, (1 - p_t)^{\gamma} \log(p_t),
\end{equation}
sweeping $\alpha$ and $\gamma$. The decision threshold is tuned post-hoc over $[0.10, 0.90]$ to enforce
recall $\geq 0.70$, after which we report precision, F1, and the threshold-free PR-AUC and ROC-AUC.

\subsection{Inductive disjoint-node evaluation}
Accounts are partitioned $80/20$ (stratified by whether an account touches any laundering edge). A
transaction is supervised in a split only if both endpoints belong to that split, so test accounts are
entirely unseen during training, mimicking the deployment setting of scoring new accounts. A $10\%$
train-internal slice forms the validation set for model selection and threshold choice.

\section{Experimental Setup}
\label{sec:setup}
\subsection{Dataset}
Table~\ref{tab:data} summarizes the data. After undersampling, the working graph has $221{,}207$ account
nodes and $247{,}247$ transaction edges; the disjoint-node split yields $158{,}829$ training edges
($6{,}336$ laundering) and $9{,}889$ test edges ($422$ laundering). The five categorical columns have
cardinalities $14, 14, 19, 19, 8$.

\begin{table}[t]
\caption{SAML-D dataset and the undersampled working graph.}
\label{tab:data}
\centering
\begin{tabular}{lr}
\toprule
Quantity & Value \\
\midrule
Full transactions & 9{,}504{,}852 \\
Laundering (full) & 9{,}873 (0.1039\%) \\
Subset edges (2.5\% majority + all positives) & 247{,}247 \\
Subset positive rate & 3.99\% \\
Account nodes & 221{,}207 \\
Train edges / laundering & 158{,}829 / 6{,}336 \\
Test edges / laundering (disjoint accounts) & 9{,}889 / 422 \\
Validation edges & 15{,}883 \\
\bottomrule
\end{tabular}
\end{table}

\subsection{Implementation}
All GraphSAGE training runs on a single Modal serverless \textbf{L4} GPU using PyTorch~2.4.1 and PyTorch
Geometric~2.6.1. We use AdamW (learning rate $10^{-3}$, weight decay $5\times10^{-4}$), gradient clipping
at $1.0$, embedding dimension $16$, a $300$-epoch cap, and early stopping with patience $20$ on validation
PR-AUC. Data is downloaded once and cached to a persistent Volume. The complete study used $15$ GPU
training runs and cost under one US dollar. GAT and Edge-GAT were excluded by design because preliminary
experiments yielded poor results, so we did not spend the budget on them.

\subsection{Baseline and metrics}
The baseline is XGBoost~\cite{chen2016xgboost} ($300$ trees, depth $6$, \texttt{scale\_pos\_weight} set to
the class ratio) on identical per-edge features. We report recall (primary, constrained $\geq 0.70$),
precision, F1, ROC-AUC, and PR-AUC, with two to three seeds per configuration and Welch $t$-tests.

\section{Results}
\subsection{Main comparison}
Table~\ref{tab:main} and Fig.~\ref{fig:main} report the headline result. GraphSAGE roughly doubles the
PR-AUC of XGBoost ($0.741$ versus $0.288$) and, at the recall $\geq 0.70$ operating point, delivers roughly
$6\times$ the precision ($0.564$ versus $0.089$), corresponding to far fewer false alerts per caught
launderer. A Welch $t$-test on PR-AUC across seeds gives $t = 76.7$, $p < 0.001$. The precision-recall
curves (Fig.~\ref{fig:pr}) show GraphSAGE dominating XGBoost across the entire operating range, and the real
SAML-D subgraph in Fig.~\ref{fig:subgraph} makes the reason concrete: laundering concentrates in dense
fan-in/fan-out hubs and repeated-transfer motifs that a per-transaction model cannot see.

\begin{table*}[t]
\caption{Main results on the disjoint-node test set (mean $\pm$ std over seeds; operating point: recall
$\geq 0.70$). Best per column in bold.}
\label{tab:main}
\centering
\begin{tabular}{lccccc}
\toprule
Model & PR-AUC & ROC-AUC & Recall & Precision & F1 \\
\midrule
XGBoost (tabular) & $0.288 \pm 0.006$ & $0.791$ & $0.742$ & $0.089 \pm 0.011$ & $0.159$ \\
GraphSAGE + weighted BCE & $\mathbf{0.741 \pm 0.013}$ & $\mathbf{0.934}$ & $0.710$ & $0.564 \pm 0.069$ & $0.627$ \\
GraphSAGE + Focal ($\alpha{=}0.5, \gamma{=}2$) & $0.736 \pm 0.008$ & $0.925$ & $0.708$ & $\mathbf{0.578 \pm 0.068}$ & $\mathbf{0.635}$ \\
\bottomrule
\end{tabular}
\end{table*}

\begin{figure*}[t]
\centering
\includegraphics[width=0.95\textwidth]{figures/output/fig3_main_results.png}
\caption{GraphSAGE versus XGBoost on SAML-D (disjoint-node test, mean $\pm$ std). Left: PR-AUC. Right:
precision at recall $\geq 0.70$. GraphSAGE roughly doubles PR-AUC and gives roughly $6\times$ the precision
of the tabular baseline at the same recall.}
\label{fig:main}
\end{figure*}

\begin{figure}[t]
\centering
\includegraphics[width=\columnwidth]{figures/output/fig_typologies_pb.png}
\caption{Money-laundering typologies as graph motifs: fan-out (smurfing), fan-in (structuring), layering
cycles, and scatter-gather. These multi-hop patterns are visible to a GNN but invisible to a per-transaction
tabular model.}
\label{fig:typologies}
\end{figure}

\begin{figure}[t]
\centering
\includegraphics[width=\columnwidth]{figures/output/fig_architecture_pb.png}
\caption{The inductive GraphSAGE encoder embeds each account from its sampled multi-hop neighborhood via
mean aggregation; the sender and receiver embeddings are concatenated with learned categorical embeddings
and numeric transaction features and passed to an edge-MLP that outputs the laundering probability.}
\label{fig:architecture}
\end{figure}

\begin{figure}[t]
\centering
\includegraphics[width=\columnwidth]{figures/output/fig1_saml_subgraph.png}
\caption{A real subgraph from SAML-D ($70$ accounts, $114$ transactions, $82$ laundering). Laundering edges
(red) concentrate in fan-in/fan-out hubs and repeated-transfer motifs; node size encodes degree.}
\label{fig:subgraph}
\end{figure}

\begin{figure}[t]
\centering
\includegraphics[width=\columnwidth]{figures/output/fig4_pr_curves.png}
\caption{Precision-recall curves on the disjoint-node test set. The recall $\geq 0.70$ target region is
shaded. GraphSAGE (both losses) dominates XGBoost across the curve.}
\label{fig:pr}
\end{figure}

\subsection{Architecture search}
Table~\ref{tab:arch} and Fig.~\ref{fig:arch} report the five-configuration architecture search under
weighted BCE. A deeper three-layer encoder is best (validation PR-AUC $0.982$, test $0.736$), consistent
with the multi-hop nature of laundering motifs. Validation PR-AUC is higher than test PR-AUC throughout
because validation edges are internal to the train graph (transductive), whereas the test set is the harder
inductive disjoint-node regime; validation is used only for selection.

\begin{table}[t]
\caption{GraphSAGE architecture search (weighted BCE, seed 0). L/H/M/d = SAGE layers / hidden / edge-MLP
layers / dropout.}
\label{tab:arch}
\centering
\begin{tabular}{llccc}
\toprule
Config & L/H/M/d & Val PR-AUC & Test PR-AUC & Prec@.70 \\
\midrule
b1 & 2/128/2/0.3 & 0.971 & 0.682 & 0.369 \\
\textbf{b2} & \textbf{3/128/2/0.3} & \textbf{0.982} & \textbf{0.736} & \textbf{0.524} \\
b3 & 2/64/2/0.3 & 0.966 & 0.685 & 0.380 \\
b4 & 2/128/3/0.3 & 0.980 & 0.709 & 0.452 \\
b5 & 2/128/2/0.5 & 0.972 & 0.695 & 0.433 \\
\bottomrule
\end{tabular}
\end{table}

\begin{figure}[t]
\centering
\includegraphics[width=\columnwidth]{figures/output/fig5_arch_search.png}
\caption{Architecture search: validation and test PR-AUC for the five GraphSAGE configurations. The star
marks the best (three-layer) encoder selected by validation PR-AUC.}
\label{fig:arch}
\end{figure}

\subsection{Focal loss versus weighted cross-entropy}
On the best architecture we sweep focal $(\alpha, \gamma)$ (Table~\ref{tab:focal}). The best setting is
$\alpha = 0.5, \gamma = 2$, notably not the vision default $\alpha = 0.25$. With three-seed confirmation,
however, focal and weighted BCE are statistically indistinguishable: PR-AUC $0.736$ versus $0.741$
($p = 0.63$) and precision at recall $\geq 0.70$ of $0.578$ versus $0.564$ ($p = 0.82$). Fig.~\ref{fig:focal}
shows the focal grid against the weighted-BCE reference band. A single-seed run had suggested a focal
advantage that multi-seed confirmation revealed to be within noise, underscoring the need for seeded error
bars in imbalanced AML evaluation.

\begin{table}[t]
\caption{Focal-loss grid on the best architecture (seed 0). n/a: no threshold reached recall $\geq 0.70$.}
\label{tab:focal}
\centering
\begin{tabular}{lccc}
\toprule
$(\alpha, \gamma)$ & Val PR-AUC & Test PR-AUC & Prec@.70 \\
\midrule
$(0.25, 1)$ & 0.964 & 0.683 & n/a \\
$(0.25, 2)$ & 0.975 & 0.714 & 0.395 \\
$(0.25, 3)$ & 0.978 & 0.729 & 0.348 \\
$(0.50, 2)$ & 0.982 & 0.744 & 0.637 \\
$(0.75, 2)$ & 0.978 & 0.742 & 0.513 \\
\bottomrule
\end{tabular}
\end{table}

\begin{figure}[t]
\centering
\includegraphics[width=\columnwidth]{figures/output/fig6_focal_vs_bce.png}
\caption{Focal-loss grid versus the weighted-BCE reference (dashed line, shaded $\pm$ std). Differences are
within seed noise: focal $\approx$ weighted BCE ($p = 0.82$ on precision at recall $\geq 0.70$).}
\label{fig:focal}
\end{figure}

\subsection{Operating point}
Fig.~\ref{fig:threshold} shows the precision/recall/F1 trade-off versus the decision threshold for the best
GraphSAGE. The recall $\geq 0.70$ constraint is met at threshold $0.4$ with precision $0.52$; lowering the
threshold trades precision for additional recall, giving practitioners an explicit, retraining-free knob to
hit a required detection rate.

\begin{figure}[t]
\centering
\includegraphics[width=\columnwidth]{figures/output/fig7_threshold.png}
\caption{Threshold sweep for the best GraphSAGE. The recall $\geq 0.70$ operating point (threshold $0.4$,
precision $0.52$) is marked.}
\label{fig:threshold}
\end{figure}

\subsection{Discussion}
Three points stand out. First, on SAML-D graph structure is the decisive lever: GraphSAGE's roughly
$6\times$ precision at fixed recall dwarfs any difference attributable to the loss function. This contrasts
with the Elliptic result where a Random Forest beat a GCN, and likely reflects SAML-D's explicit typological
graph structure, which an inductive encoder captures well. Second, focal loss does not robustly beat
weighted cross-entropy here; once majority undersampling and threshold tuning are in place, the loss choice
is secondary. Third, a credible AML-GNN benchmark does not require large-scale compute: the entire study fit
in under one dollar on a single L4 GPU.

\FloatBarrier
\section{Limitations and Future Work}
SAML-D is synthetic and its typologies may favor structural models; results should be validated on real or
additional benchmarks (AMLworld, Elliptic). We study one GNN family (GraphSAGE-mean) and a $2.5\%$
undersample rather than the full graph, and we do not model temporal dynamics. Future work includes
real-data validation, dynamic and advanced-aggregation GNNs, full-graph mini-batch training, and automatic
per-layer rank or neighbor-budget allocation under the same recall-first, significance-tested protocol.

\section{Conclusion}
On the SAML-D benchmark, an inductive GraphSAGE edge classifier decisively outperforms a strong XGBoost
baseline for money-laundering detection (PR-AUC $0.74$ versus $0.29$, roughly $6\times$ precision at recall
$\geq 0.70$, $p < 0.001$), while focal loss is statistically on par with weighted cross-entropy. Relational
structure, not the loss function, is the decisive lever, and a reproducible benchmark fits under one dollar
on a single L4 GPU.

\section*{Acknowledgments}
The authors thank Vizuara AI Labs for providing mentorship and computational resources for this research.

\begin{thebibliography}{99}
\bibitem{hamilton2017graphsage} W. L. Hamilton, R. Ying, and J. Leskovec, ``Inductive representation learning on large graphs,'' in \textit{Proc. NeurIPS}, 2017. arXiv:1706.02216.
\bibitem{oztas2023samld} B. Oztas, D. Cetinkaya, F. Adedoyin, M. Budka, H. Dogan, and G. Aksu, ``Enhancing anti-money laundering: development of a synthetic transaction monitoring dataset,'' in \textit{Proc. IEEE Int. Conf. e-Business Engineering (ICEBE)}, 2023, pp. 47--54.
\bibitem{lin2017focal} T.-Y. Lin, P. Goyal, R. Girshick, K. He, and P. Doll\'ar, ``Focal loss for dense object detection,'' in \textit{Proc. ICCV}, 2017. arXiv:1708.02002.
\bibitem{weber2019aml} M. Weber, G. Domeniconi, J. Chen, D. K. I. Weidele, C. Bellei, T. Robinson, and C. E. Leiserson, ``Anti-money laundering in Bitcoin: experimenting with graph convolutional networks for financial forensics,'' in \textit{KDD Workshop on Anomaly Detection in Finance}, 2019. arXiv:1908.02591.
\bibitem{altman2023amlworld} E. Altman, J. Blanu\v{s}a, L. von Niederh\"ausern, B. Egressy, A. Anghel, and K. Atasu, ``Realistic synthetic financial transactions for anti-money laundering models,'' in \textit{Proc. NeurIPS Datasets and Benchmarks}, 2023. arXiv:2306.16424.
\bibitem{dou2020caregnn} Y. Dou, Z. Liu, L. Sun, Y. Deng, H. Peng, and P. S. Yu, ``Enhancing graph neural network-based fraud detectors against camouflaged fraudsters,'' in \textit{Proc. CIKM}, 2020. arXiv:2008.08692.
\bibitem{liu2021pcgnn} Y. Liu, X. Ao, Z. Qin, J. Chi, J. Feng, H. Yang, and Q. He, ``Pick and choose: a GNN-based imbalanced learning approach for fraud detection,'' in \textit{Proc. The Web Conf. (WWW)}, 2021.
\bibitem{chen2016xgboost} T. Chen and C. Guestrin, ``XGBoost: a scalable tree boosting system,'' in \textit{Proc. KDD}, 2016, pp. 785--794. arXiv:1603.02754.
\bibitem{cheng2024review} D. Cheng, Y. Zou, S. Xiang, and C. Jiang, ``Graph neural networks for financial fraud detection: a review,'' \textit{Frontiers of Computer Science}, 2024. arXiv:2411.05815.
\bibitem{velickovic2018gat} P. Veli\v{c}kovi\'c, G. Cucurull, A. Casanova, A. Romero, P. Li\`o, and Y. Bengio, ``Graph attention networks,'' in \textit{Proc. ICLR}, 2018. arXiv:1710.10903.
\bibitem{kipf2017gcn} T. N. Kipf and M. Welling, ``Semi-supervised classification with graph convolutional networks,'' in \textit{Proc. ICLR}, 2017. arXiv:1609.02907.
\bibitem{davis2006pr} J. Davis and M. Goadrich, ``The relationship between precision-recall and ROC curves,'' in \textit{Proc. ICML}, 2006.
\bibitem{fey2019pyg} M. Fey and J. E. Lenssen, ``Fast graph representation learning with PyTorch Geometric,'' in \textit{ICLR Workshop on Representation Learning on Graphs and Manifolds}, 2019. arXiv:1903.02428.
\bibitem{chawla2002smote} N. V. Chawla, K. W. Bowyer, L. O. Hall, and W. P. Kegelmeyer, ``SMOTE: synthetic minority over-sampling technique,'' \textit{J. Artificial Intelligence Research}, vol. 16, pp. 321--357, 2002.
\bibitem{dalpozzolo2015} A. Dal Pozzolo, O. Caelen, R. A. Johnson, and G. Bontempi, ``Calibrating probability with undersampling for unbalanced classification,'' in \textit{Proc. IEEE SSCI}, 2015.
\bibitem{pareja2020evolvegcn} A. Pareja et al., ``EvolveGCN: evolving graph convolutional networks for dynamic graphs,'' in \textit{Proc. AAAI}, 2020. arXiv:1902.10191.
\bibitem{lopezrojas2016paysim} E. A. Lopez-Rojas, A. Elmir, and S. Axelsson, ``PaySim: a financial mobile money simulator for fraud detection,'' in \textit{Proc. EMSS}, 2016.
\bibitem{suzumura2021amlsim} T. Suzumura and H. Kanezashi, ``AMLSim: anti-money laundering simulator,'' IBM, 2021.
\bibitem{oztas2022} B. Oztas, D. Cetinkaya, F. Adedoyin, and M. Budka, ``Enhancing transaction monitoring controls to detect money laundering using machine learning,'' in \textit{Proc. IEEE ICEBE}, 2022.
\bibitem{hamilton2020book} W. L. Hamilton, \textit{Graph Representation Learning}. Morgan \& Claypool, 2020.
\end{thebibliography}

\end{document}
